% Keep the horizon evidence immediately beside its discussion.
\begingroup
\setlength{\intextsep}{8pt}
\begin{table}[H]
\caption{\PaperRevision{前瞻范围对编译结果的影响}}
\label{tab:sensitivity}
\centering
\PaperTableSetup
\fontsize{8.5pt}{9.7pt}\selectfont
\setlength{\tabcolsep}{2pt}
\begin{tabular*}{\columnwidth}{@{}l@{\extracolsep{\fill}}rrrrr@{}}
\toprule
前瞻层数上限 $H$ & 0 & 1 & 2 & 4 & 8\\
\midrule
\PaperRevision{保真度比} & \HorizonExtZeroFidelityRatio & \HorizonExtOneFidelityRatio & \HorizonExtTwoFidelityRatio & \HorizonExtFourFidelityRatio & \HorizonExtEightFidelityRatio\\
\PaperRevision{重排批次比} & \HorizonExtZeroBatchRatio & \HorizonExtOneBatchRatio & \HorizonExtTwoBatchRatio & \HorizonExtFourBatchRatio & \HorizonExtEightBatchRatio\\
\bottomrule
\end{tabular*}
\PaperTableNote{\fontsize{8pt}{9.1pt}\selectfont\PaperRevision{注：各项相对同一电路的$H=8$结果，参照为1；$H=0$仅评价当前层。\HorizonExtTargetN{}个电路中，\HorizonExtCompleteN{}个在五档配置下均满足比较条件。逐电路、逐档取\HorizonExtSeedN{}次结果的中位数，与该电路$H=8$的中位数作比，再取几何均值。结果保留四位小数。}}
\end{table}
\endgroup
